\chapter{Type Analysis}
\label{ch:typ}

% One of the main challenges of the analysis of Rego policies is its dynamic type system, which contrasts the static typing of modern SMT solvers.
% To create correct models of the policies, we 
%
To effectively and correctly create a SMT model of policy specified in Rego language, it is important to infer datatypes of appearing values, as precisely as possible.
In this chapter, I describe a set of rules followed by the type analysis, which creates a function $\TA$, mapping every value to its corresponding type.
This chapter will be heavily reliant on the definition of Rego type system from \ref{s:typ}.

\section{Properties of Types}

This section introduces additional properties of Rego types, that it their \emph{depth} and \emph{precision order}.
Additionally, this section presents \emph{union normalization}, i.e., a formal process to transform the union type into a standard (normalized) form.

\subsection{Type depth}

The \emph{depth} $\depth$ of a types is defined as follows:
\begin{itemize}
  \item $\depth(\ty) = 0$, for $\ty$ being an atomic type,
  \item $\depth(\ty_1 \cup \cdots \cup \ty_m) = \depth(\ty_i)$, for $1 \leq i \leq m$, where $\forall_{1 \leq j \leq m}:\depth(\ty_i) \geq \depth(\ty_j)$,
  \item $\depth(\Obj{f_1{:}\ty_1,\ldots,f_n{:}\ty_n}) = 1+\depth(\ty_1 \cup \cdots \cup \ty_n)$,
  \item $\depth(\Array{\ty}) = 1+\depth(\ty)$,
  \item $\depth(\RegoSet{\ty}) = 1+\depth(\ty)$.
\end{itemize}

Informally said, atomic types have depth $0$, and composite types have their depth greater by one than their element with the maximum depth.
In addition, we allow for the type depth function $\depth$ to be applied directly to terms with the intuitive semantics: $\depth(t) = \depth(\typof(t))$, where $t \in \Terms$.

% % TODO
% \subsection{Precision Order}
%
% Precision order of types if defined as the smallest preorder satisfying:
% \begin{enumerate}
%   \item $\Unknown \leqT \ty$ for all $\ty \in \mathbb{T}$.
%   \item $\ty \leqT \ty$ (reflexivity).
%   \item $\Array{\ty} \leqT \Array{\ty'}$ iff $\ty \leqT \ty'$.
%   \item $\Obj{f_i{:}\ty_i} \leqT \Obj{g_j{:}\ty_j'}$
%         iff $\{g_j\} \subseteq \{f_i\}$
%         and $\ty_{g_j} \leqT \ty_{g_j}'$ for all $j$.
%         (More fields, or more precise field types, is more precise.)
%   \item For unions: $(\ty_1 \cup \cdots \cup \ty_m) \leqT (\ty_1' \cup \cdots \cup \ty_n')$ iff $\{\ty_i\} \subseteq \{\ty_j'\}$ (the set of members of the left union is contained in the set of members of the right).
% \end{enumerate}
% We write $\ty \ltT \ty'$ for $\ty \leqT \ty'$ and $\ty' \not\leqT \ty$ (strict precision).
%
% The precision order formalises the intuition that \emph{more information is better}: $\ty \leqT \ty'$ reads ``$\ty$ is at most as informative as $\ty'$''.
% Unknown is the least informative type.
% An object with more fields is strictly more informative than one with fewer fields.
% A union with fewer alternatives is more informative than one with more.

% \subsection{Union Canonicalization}
%
% Union types are subject to a \emph{canonicalization} operation that keeps them in a normal form.
%
% \begin{definition}[Canonicalization $\canon$]
% \label{def:canon}
% Given a type $\ty$, $\canon(\ty)$ is computed as follows:
% \begin{enumerate}
%   \item If $\ty$ is not a union, return $\ty$ unchanged.
%   \item \emph{Flatten}: replace every union member that is itself a union by the transitive closure of its members.
%   \item \emph{Deduplicate}: remove equal members (by structural equality), keeping first occurrences.
%   \item \emph{Drop unknowns}: if the resulting set has more than one member, remove all $\Unknown$ members.
%   \item \emph{Simplify}: if exactly one member remains, return that member (not a union).
% \end{enumerate}
% \end{definition}
%
% The effect of $\canon$ is that the empty union is the unique representation for ``no information'', and single-member unions are collapsed.
% After canonicalization, a union $\ty_1 \cup \cdots \cup \ty_m$ has $m \geq 2$ and no $\Unknown$ member (unless it is the sole member of the whole type, which would have been simplified away).
%
% \section{Typing Environments}
% \label{sec:env}
% %%=============================================================
%
% The type analyzer is parameterised by a \emph{typing context}
% \[
% \Sigma \;=\; (\typof,\; \schema,\; \params,\; \pkg)
% \]
% where:
% \begin{itemize}
%   \item $\typof : \mathsf{Key} \to \mathbb{T}$ is a mutable \emph{type map} from string keys (stringified AST values) to types.
%     We write $\typof(v)$ for $\typof(\mathsf{key}(v))$.
%   \item $\schema : \mathsf{Path} \rightharpoonup \mathbb{T}$ is an \emph{input schema} providing types for paths into \texttt{input.review.object}.
%     It is built from an example YAML/JSON document or a JSON Schema.
%   \item $\params : \mathsf{Name} \rightharpoonup \mathbb{T}$ is a \emph{parameter map} for \texttt{input.parameters.\textit{name}} references.
%   \item $\pkg$ is the \emph{package reference} (a sequence of names) used to resolve \texttt{data.\textit{pkg}.\textit{rule}} references.
% \end{itemize}
%
% The type map $\typof$ is updated in place during analysis.
% Two update operations are used:
%
% \paragraph{Precision-based update (\texttt{setType}).}
% $\typof[v \mapsto \ty]$ overwrites the existing entry for $v$ only if $\typof(v) \leqT \ty$ (strictly less precise):
% \[
% \typof \mathrel{[v \xrightarrow{\leqT} \ty]} =
% \begin{cases}
% \typof[v \mapsto \ty] & \text{if } \ty \ltT \typof(v) \\
% \typof & \text{otherwise}
% \end{cases}
% \]
%
% \paragraph{Union extension (\texttt{addToType}).}
% $\typof \mathrel{[v \xrightarrow{\cup} \ty]}$ converts the current type of $v$ into a union and appends $\ty$, then canonicalizes:
% \[
% \typof \mathrel{[v \xrightarrow{\cup} \ty]} = \typof\bigl[v \mapsto \canon(\typof(v) \cup \ty)\bigr]
% \]
%
% %%=============================================================
% \section{Type Inference Rules}
% \label{sec:inference}
% %%=============================================================
%
% We use three judgment forms:
% \begin{align*}
% \Sigma \vdashT v : \ty  &\quad\text{(term/value type inference)} \\
% \Sigma \vdashE e : \ty  &\quad\text{(expression type inference)} \\
% \Sigma \vdashR \rho : \ty &\quad\text{(reference type inference)}
% \end{align*}
% All judgments have side effects on $\typof$.
% The notation $\Sigma \vdashT v : \ty \leadsto \Sigma'$ makes the update explicit when needed; we omit it when the state change is clear from context.
%
% \subsection{Term Inference (\texttt{inferAstType})}
% \label{sec:term-inference}
%
% \paragraph{Caching.}
% Before case analysis, the system checks the cache:
% if $v \in \dom(\typof)$ and $\typof(v) \neq \Unknown$, the cached type is returned directly.
%
% % \begin{figure}[h]
% % \[
% % \textbf{(T-Nil)}\quad
% % \inferrule{ }{\Sigma \vdashT \mathbf{nil} : \Unknown}
% % \qquad
% % \textbf{(T-String)}\quad
% % \inferrule{ }{\Sigma \vdashT s : \Str}
% % \qquad
% % \textbf{(T-Number)}\quad
% % \inferrule{ }{\Sigma \vdashT n : \Int}
% % \]
% %
% % \[
% % \textbf{(T-Bool)}\quad
% % \inferrule{ }{\Sigma \vdashT b : \Bool}
% % \qquad
% % \textbf{(T-Set)}\quad
% % \inferrule{ }{\Sigma \vdashT \{t_1,\ldots\} : \RegoSet}
% % \]
%
% \[
% \textbf{(T-ArrayEmpty)}\quad
% \inferrule{ }{\Sigma \vdashT [] : \Array{\Unknown}}
% \qquad
% \textbf{(T-ArrayNonEmpty)}\quad
% \inferrule{
%   \Sigma \vdashT v_0 : \ty
% }{
%   \Sigma \vdashT [v_0, v_1, \ldots] : \Array{\ty}
% }
% \]
%
% \[
% \textbf{(T-Object)}\quad
% \inferrule{
%   \Sigma \vdashT v_i : \ty_i \quad \forall i \in 1..n
% }{
%   \Sigma \vdashT \{``k_1"{:}v_1, \ldots, ``k_n"{:}v_n\} : \Obj{k_1{:}\ty_1, \ldots, k_n{:}\ty_n}
% }
% \]
% \caption{Type inference rules for literal terms.}
% \label{fig:literal-rules}
% \end{figure}
%
% \begin{remark}
% Rule \textbf{(T-ArrayNonEmpty)} infers the element type from the \emph{first} element only.
% This reflects the homogeneous-array assumption: in policy code, arrays are expected to be uniform.
% Heterogeneous arrays produce a union type via the union-of-elements generalisation, but the current implementation chooses a simpler over-approximation for efficiency.
% \end{remark}
%
% \paragraph{Variable inference.}
% \[
% \textbf{(T-VarLookup)}\quad
% \inferrule{
%   x \in \dom(\typof)
% }{
%   \Sigma \vdashT x : \typof(x)
% }
% \qquad
% \textbf{(T-VarUnknown)}\quad
% \inferrule{
%   x \notin \dom(\typof)
% }{
%   \Sigma \vdashT x : \Unknown
% }
% \]
%
% Type propagation from a function parameter expectation:
% \[
% \textbf{(T-VarRefine)}\quad
% \inferrule{
%   \Sigma \vdashT x : \ty \\
%   \ty_{\mathrm{exp}} \ltT \ty
% }{
%   \Sigma \vdashT x : \ty_{\mathrm{exp}}
% }
% \]
% Rule \textbf{(T-VarRefine)} applies when an \emph{inherited type} $\ty_{\mathrm{exp}}$ (from a function parameter expectation) is strictly more precise than the currently known type $\ty$: in that case the expected type takes priority.
%
% \paragraph{Reference inference.}
% \[
% \textbf{(T-Ref)}\quad
% \inferrule{
%   \Sigma \vdashR r : \ty
% }{
%   \Sigma \vdashT r : \ty
% }
% \]
%
% \paragraph{Call-term inference.}
% \[
% \textbf{(T-Call)}\quad
% \inferrule{
%   \mathsf{sig}(f, k) = (\ty_r, \ty_1, \ldots, \ty_k) \\
%   \Sigma \vdashT t_i : \ty_i' \quad \forall i \in 1..k
% }{
%   \Sigma \vdashT f(t_1,\ldots,t_k) : \ty_r
% }
% \]
% where $\mathsf{sig}(f,k)$ is the predefined signature for function $f$ with $k$ parameters, providing the return type $\ty_r$ and the expected parameter types $\ty_1,\ldots,\ty_k$ used as $\ty_{\mathrm{exp}}$ when inferring each argument.
%
% \subsection{Expression Inference (\texttt{InferExprType})}
% \label{sec:expr-inference}
%
% \begin{figure}[h]
% \[
% \textbf{(E-Term)}\quad
% \inferrule{
%   \Sigma \vdashT t : \ty
% }{
%   \Sigma \vdashE t : \ty
% }
% \]
%
% \[
% \textbf{(E-Eq)}\quad
% \inferrule{
%   \Sigma \vdashT t_L : \ty_L \\
%   \Sigma \vdashT t_R : \ty_R \\
%   t_L \in \mathsf{Var} \Rightarrow \typof \mathrel{[t_L \xrightarrow{\cup} \ty_R]} \\
%   t_R \in \mathsf{Var} \Rightarrow \typof \mathrel{[t_R \xrightarrow{\cup} \ty_L]}
% }{
%   \Sigma \vdashE (t_L \mathrel{=} t_R) : \Bool
% }
% \]
%
% \[
% \textbf{(E-Op)}\quad
% \inferrule{
%   \mathit{op} \in \{<, \leq, >, \geq, \neq\}
% }{
%   \Sigma \vdashE (t_L \;\mathit{op}\; t_R) : \Bool
% }
% \qquad
% \textbf{(E-Call)}\quad
% \inferrule{
%   \mathsf{sig}(f, k) = (\ty_r, \ty_1, \ldots, \ty_k) \\
%   \Sigma \vdashT t_i : \_ \quad \forall i
% }{
%   \Sigma \vdashE f(t_1,\ldots,t_k) : \ty_r
% }
% \]
% \caption{Type inference rules for expressions.}
% \label{fig:expr-rules}
% \end{figure}
%
% Rule \textbf{(E-Eq)} is the primary mechanism for bidirectional type propagation between variables.
% When a variable $x$ is equated with a term $t$ of type $\ty$, the type $\ty$ is \emph{unioned into} the type of $x$ (and vice versa if both sides are variables).
% This reflects Rego's unification semantics: the equality statement does not simply assign but \emph{constrains}, and both sides carry information about each other.
%
% \subsection{Reference Inference (\texttt{inferRefType})}
% \label{sec:ref-inference}
%
% Let $r = v[s_1]\cdots[s_k]$ where $v$ is the head term.
%
% \[
% \textbf{(R-Empty)}\quad
% \inferrule{ }{\Sigma \vdashR \langle\;\rangle : \Unknown}
% \]
%
% \[
% \textbf{(R-InputParam)}\quad
% \inferrule{
%   r = \texttt{input.parameters.}n \\
%   n \in \dom(\params)
% }{
%   \Sigma \vdashR r : \params(n)
% }
% \]
%
% \[
% \textbf{(R-InputSchema)}\quad
% \inferrule{
%   r = \texttt{input.review.object.}p \\
%   \schema(p) = \ty
% }{
%   \Sigma \vdashR r : \ty
% }
% \]
%
% \[
% \textbf{(R-Data)}\quad
% \inferrule{
%   r = \texttt{data.}\pkg.q.p' \\
%   q \in \dom(\typof) \\
%   \pathtype(\typof(q), p') = \ty
% }{
%   \Sigma \vdashR r : \ty
% }
% \]
%
% \[
% \textbf{(R-VarHead)}\quad
% \inferrule{
%   r = x[s_1]\cdots[s_k] \\
%   x \in \dom(\typof) \\
%   \pathtype(\typof(x), [s_1,\ldots,s_k]) = \ty
% }{
%   \Sigma \vdashR r : \ty
% }
% \qquad
% \textbf{(R-Unknown)}\quad
% \inferrule{
%   \text{(no other rule applies)}
% }{
%   \Sigma \vdashR r : \Unknown
% }
% \]
%
% The auxiliary function $\pathtype(\ty, p)$ traverses type $\ty$ along path $p$ and is defined in \S\ref{sec:path}.
%
% \subsection{Path Traversal (\texttt{GetTypeFromPath})}
% \label{sec:path}
%
% Given a type $\ty$ and a sequence of path nodes $p = [n_1, \ldots, n_\ell]$, the function $\pathtype(\ty, p)$ resolves the nested type:
%
% \[
% \textbf{(P-Empty)}\quad
% \inferrule{ }{\pathtype(\ty, []) = \ty}
% \]
%
% \[
% \textbf{(P-ObjGround)}\quad
% \inferrule{
%   \ty = \Obj{\ldots, f{:}\ty_f, \ldots} \\
%   n_1 = f \text{ (ground key)}
% }{
%   \pathtype(\ty, [f, n_2, \ldots]) = \pathtype(\ty_f, [n_2,\ldots])
% }
% \]
%
% \[
% \textbf{(P-ObjNonGround)}\quad
% \inferrule{
%   \ty = \Obj{f_1{:}\ty_1, \ldots, f_m{:}\ty_m} \\
%   n_1 \text{ is a variable (non-ground)} \\
%   \ty_\cup = \canon(\ty_1 \cup \cdots \cup \ty_m)
% }{
%   \pathtype(\ty, [n_1, n_2, \ldots]) = \pathtype(\ty_\cup, [n_2, \ldots])
% }
% \]
%
% \[
% \textbf{(P-Array)}\quad
% \inferrule{
%   \ty = \Array{\ty'}
% }{
%   \pathtype(\ty, [n_1, n_2, \ldots]) = \pathtype(\ty', [n_2, \ldots])
% }
% \]
%
% \[
% \textbf{(P-Union)}\quad
% \inferrule{
%   \ty = \ty_1 \cup \cdots \cup \ty_m \\
%   \pathtype(\ty_i, p) = \ty_i' \quad \forall i
% }{
%   \pathtype(\ty, p) = \canon(\ty_1' \cup \cdots \cup \ty_m')
% }
% \]
%
% Rule \textbf{(P-ObjNonGround)} handles the common Rego idiom \texttt{x := obj[\_]}: the wildcard \texttt{\_} denotes an unspecified key, so the type of $x$ is the union of all field types of the object.
%
